\documentclass[12pt]{letter}
\usepackage[margin=2.5cm]{geometry}
\usepackage[T1]{fontenc}
\usepackage[utf8]{inputenc}
\usepackage{hyperref}

\signature{David Alfyorov}
\address{Independent Researcher\\davidich.alfyorov@gmail.com}

\begin{document}
\begin{letter}{Editorial Board\\General Relativity and Gravitation}

\opening{Dear Editor,}

Please consider the enclosed manuscript, ``The predictive content
of Spectral Causal Theory,'' for publication as a Research Article
in \textit{General Relativity and Gravitation}.

Spectral Causal Theory (SCT) is a one-loop gravitational effective
field theory built from the spectral action of a noncommutative
geometry spectral triple encoding the Standard Model.  Seven papers
have reported results from this program; the present paper collects
the predictions that follow and classifies them by their dependence
on the cutoff function.

The paper contains original computation that has not appeared
elsewhere: a scan of propagator zeros across the cutoff family
$f(u)=e^{-u^n}$ ($n=1,\dots,5$), a derivation of the logarithmic
black hole entropy coefficient $c_{\log}=37/24$ from the Sen
formula, a comparison with five quantum gravity programs (LQG, AS,
CDT, string theory, IDG) across nine quantitative axes, and five
falsification criteria tied to forthcoming experiments.  All
underlying results have been cross-checked numerically to 100-digit
precision, verified in Lean~4, and confirmed by independent
re-derivation.

I believe the manuscript fits the scope of GRG: it concerns the
predictive and testable content of a quantum gravity framework and
provides a quantitative comparison with competing approaches.

\closing{Respectfully,}

\end{letter}
\end{document}
